\scanpage{85}
\mspage{\handwritten{69}}\mspage{\handwritten{1}}

% A single diagonal stroke runs through the whole typescript.
\deleted{Appliquons \deleted{\unsure{(1|1.|1}}{\handwritten{7}}.2.1\deleted{\unsure{)|à}} au groupe formel $\handwritten{\hat}{A}$, dans le cas favorable envisagé dans \deleted{\unsure{le|7.0|la}} 7.1 où le schéma \deleted{\unsure{abél|Abél|ab.}} formel abélien $\handwritten{\hat}{B}$ est algébrisable\handwritten{,} \inserted{\handwritten{et où on suppose de plus \deleted{\unsure{le}} $T_0$ isotrivial.}} On trouve alors un\deleted{e extension canonique du} schéma \deleted{abélien $B$ dont le}\inserted{\handwritten{en groupes $A^{\unsure{\natural|\sharp}}$ sur $S$}}, \inserted{\handwritten{appelé le \uline{groupe de Raynaud} associé au schéma en groupes $A$ sur $S$\unsure{,|$\vphantom{S}_{o}$,} tel que $A$}} complété formel est $B$, par le tore $T$ sur $S$ qui remonte $T_0$\handwritten{,}\deleted{\unsure{;|,|et}} que nous ap\handwritten{p}ellerons l'\uline{extension de Raynaud} associée à $A$, et que nous noterons $A^{\handwritten{\natural}}$ :
\[
 0\handwritten{\longrightarrow}T\handwritten{\longrightarrow}A^{\handwritten{\natural}}\handwritten{\longrightarrow}B\handwritten{\longrightarrow}0\ .
 \tag*{(\deleted{1}\handwritten{7}.2.2)}
\]
Elle est donc définie, à isomorphisme unique près, par le fait qu'elle donne naissance à un système cohérent d'isomorphismes d'extensions
\[
 A^{\handwritten{\natural}}_n\handwritten{\simeq}A_n
 \tag*{(7.2.\handwritten{3})}
\]
(le deuxième membre muni de sa structure d'extension décrite dans (7.1.1) i.e. qu'on a un isomorphisme de \inserted{schémas formels en} groupes \deleted{formels}
\[
 (A^{\handwritten{\natural}})^{\handwritten{\wedge}}\handwritten{\simeq}\handwritten{\hat}{A}\ ,
 \tag*{(7.\deleted{3}\handwritten{2}.4)}
\]
respectant d'ailleurs les structures d'extensions. \deleted{Il résulte de plus de 7.2.1}

7.2.5. Il résulte d'autre part \deleted{que la} immédiatement de 7.2.1 que la suite exacte (7.2.2) dépend fonctoriellement du schéma \inserted{en groupes} $A$ sur $S$ satisfaisant aux conditions envisagées dans 7.1, et que sa formation commute à tout changement de base par un homomorphisme $S'\handwritten{\to} S$ associé à un homomorphisme d'anneaux topologiques noethériens adiques $V\handwritten{\to}V'$.}
% archive sheet number: 84
